
\subsubsection{LLM Calibration and Reliability Under Alignment}

Liu et al. [2023] demonstrate that instruction fine-tuning alters calibration across a range of tasks, with Expected Calibration Error increasing even when accuracy improves. Li et al. [2024] document that RLHF produces heterogeneous effects across the input distribution rather than uniform shifts: different inputs exhibit qualitatively different responses to the same alignment procedure. Xu et al. [2024] compare PPO and DPO algorithmically and show that the two alignment paradigms produce systematically different behavioral profiles despite optimizing related preference objectives. Wang et al. [2024] provide a comprehensive survey of alignment techniques including RLHF, RLAIF, PPO, and DPO. This literature establishes that alignment is a structured transformation of model behavior with input-dependent effects.

\subsubsection{Pre-Alignment Predictors of Post-Alignment Behavior}

Plaut et al. [2024] demonstrate that the maximum softmax probability of a base model predicts whether the model's answer is correct, establishing that pre-alignment confidence carries predictive signal about answer quality. Fan et al. [2026] examine correlations between pre-alignment benchmark accuracy and post-alignment accuracy, finding that aggregate accuracy trends are partially predictable from base model performance. These works operate at the aggregate level --- across items or benchmarks --- rather than predicting per-item change. The present work differs in target: it asks not whether the base model answer is correct, but whether alignment will change it, and operates at the individual item level.

\subsubsection{Geometric Analysis of Model Representations}

Khanmohammadi et al. [2025] investigate how representation stability relates to model calibration, finding that geometric properties of the representation space carry information about model reliability. Lunardi et al. [2025] examine MCQ answer reliability under paraphrase perturbations, finding that model responses are more fragile than surface accuracy suggests, a finding they attribute partly to proximity of items to decision boundaries in the representation space. This geometric perspective motivates the non-isotropy analysis conducted in this paper.

\subsubsection{Gap and Positioning}

The literature reviewed above establishes that alignment reshapes model calibration with heterogeneous effects, that base model signals predict aggregate post-alignment properties, and that geometric properties of representations are informative about reliability. No prior work poses the per-item, pre-deployment prediction of argmax instability as a classification task, operationalizes it with the confidence margin, or evaluates it with cross-benchmark generalization. The present paper addresses this gap.

